\begin{textsample}{POS Dim 5 – global\_north\_2024\_02 – Score 15.00 – global\_north\_2024\_02/106036111.txt}  \label{ex:f5_pos_250}
Palestinian Foreign Minister Riyad al-Maliki demanded an immediate end to Israel ' s occupation of the Palestinian territories at the start of hearings on Monday into the \textbf{legal} status of the occupation at the United Nations ' top \textbf{court}.
More than 50 states will present \textbf{arguments} before the \textbf{International} \textbf{Court} of Justice ( ICJ ) in The Hague until Feb. 26, following a 2022 request from the U.N.
General Assembly for an \textbf{advisory}, or \textbf{non-binding}, opinion.
Al-Maliki accused Israel of subjecting Palestinians to decades of discrimination and apartheid - \textbf{accusations} Israel has rejected - arguing that they had been left with the choice of"displacement, subjugation, or death"."The only solution consistent with \textbf{international} \textbf{law} is for this illegal occupation to come to an immediate, \textbf{unconditional} and total end,"he said.
The ICJ ' s 15-judge panel has been asked to review Israel ' s"occupation, settlement and annexation... including \textbf{measures} aimed at altering the demographic composition, character and status of the Holy City of Jerusalem, and from its adoption of related discriminatory legislation and \textbf{measures} @ @ @ @ @ @ @ @ @ @ six months to issue an opinion on the request, which also asks them to consider the \textbf{legal} status of the occupation and its consequences.
Israel captured the West Bank, Gaza and East Jerusalem - areas of historic Palestine which the Palestinians want for a state - in a 1967 war and has since built settlements in the West Bank and steadily expanded them.
Israeli leaders have long disputed that the territories are formally occupied on the basis that they were captured from Jordan and Egypt during a war rather than from a sovereign Palestine.
The United Nations has referred to the territories as occupied by Israel since 1967 and demanded that Israeli forces withdraw, saying it is the only way to secure peace.
Its 1967 \textbf{resolution} did not, however, specifically label the \textbf{resolution} as illegal."The best and possibly the last hope for the two-state solution, so vital to the needs of both peoples, is for the \textbf{court} to declare illegal the main obstacle to that solution : the ongoing Israeli occupation of @ @ @ @ @ @ @ @ @ @ Palestinians, told \textbf{judges}.
Israel will not attend the hearings but has sent written observations.
While Israel has ignored \textbf{legal} opinions in the past, this one could increase political pressure over its war in Gaza, which has killed about 29,000 Palestinians, according to Gaza health officials, since Hamas attacked Israel on Oct. 7.
It withdrew from Gaza in 2005, but, along with neighbouring Egypt, still controls its borders.
It has also annexed East Jerusalem in a move not \textbf{recognised} by most countries. ' MORAL, POLITICAL AND \textbf{LEGAL} IMPERATIVE ' The hearing is part of Palestinian efforts to get \textbf{international} \textbf{legal} institutions to examine Israel ' s conduct.
These have stepped up since Israel ' s war on Gaza in response to the Hamas attacks, which killed 1,200 people, according to Israeli tallies.
Al-Maliki reiterated \textbf{accusations} of Israeli \textbf{genocide} in Gaza which Israel rebuffed at separate hearings in The Hague last month in which the World \textbf{Court} ordered Israel to do everything in its power to prevent \textbf{acts} of \textbf{genocide} @ @ @ @ @ @ @ @ @ @ result of decades of impunity and inaction.
Ending Israel ' s impunity is a moral, political and \textbf{legal} imperative,"al-Maliki said.
Israel has said it faces an existential threat by Hamas militants and other groups and is \textbf{acting} in self-defence.
There are mounting concerns about an Israeli ground offensive against the Gaza city of Rafah, a last refuge for more than a million Palestinians after they fled to the south of the enclave to avoid Israeli assaults.
It is the second time the U.N.
General Assembly has asked the ICJ, also known as the World \textbf{Court}, for an \textbf{advisory} opinion related to the occupied Palestinian territory.
In July 2004, the \textbf{court} found that Israel ' s separation wall in the West Bank violated \textbf{international} \textbf{law} and should be dismantled, though it still stands to this day.
Disclaimer : The content of this article is syndicated or provided to this website from an external third party provider.
We are not responsible for, and do not control, such external websites, @ @ @ @ @ @ @ @ @ @ the text is provided on an"as is"and"as available"basis and has not been edited in any way.
Neither we nor our affiliates guarantee the accuracy of or endorse the views or opinions expressed in this article.
Read our full disclaimer policy here.

% matched lemmas: accusation, act, advisory, argument, court, genocide, international, judge, law, legal, measure, non-binding, recognise, resolution, unconditional
\end{textsample}
